% AUTORIDAD DE FUENTE: Exposé IX §§12--12.2, líneas francesas
% 3452--3547; índices de escaneo 476--480, folios 465--469.

%%%%%%%%%%  scan idx 476  |  folio 465  %%%%%%%%%%

\medskip
\noindent
12. \underline{Modelo de Néron de una jacobiana y fórmula de
Picard--Lefschetz}

Recordemos primero los siguientes resultados generales de M. Raynaud
[22, teoremas 5 y 7; 22 bis].

\medskip
\noindent
\underline{Teorema} 12.1 (M. Raynaud). \underline{Sean $S$ un trait
estrictamente local y $f:X\to S$ un esquema propio y plano sobre $S$, de
fibras de dimensión uno, tales que}

\begin{equation}\tag{12.1.1}
 f_*\underline{\mathcal O}_X\simeq\underline{\mathcal O}_S,
\end{equation}

\noindent
\underline{y tales que los anillos locales de $X$ sean factoriales. Para
cada punto maximal $x_i$ de $X_o$, pongamos}

\begin{equation}\tag{12.1.2}
 d_i=\operatorname{long}\underline{\mathcal O}_{X_o,x_i},
 \qquad d_i^t=\mu_i d_i,
\end{equation}

\noindent
\underline{donde $\mu_i$ es la multiplicidad radicial de $k(x_i)$ sobre
$k$} (EGA IV, 4.7.4) (\underline{por tanto, una potencia del exponente
característico de} $k$). \underline{Sea $d$ (respectivamente, $d^t$) el máximo común
divisor de los enteros $d_i$ (respectivamente, $d_i^t$); se tiene, pues,
$d=d^t$ si $k$ es perfecto. Pongamos}

\begin{equation}\tag{12.1.3}
 P=\underline{\operatorname{Pic}}_{X/S}:
   (\operatorname{Sch})_{/S}^{\,o}\longrightarrow(\operatorname{Ab}).
\end{equation}

\noindent
a) \underline{Consideremos las condiciones siguientes}:

\smallskip
\noindent
1) $d^t=1$.

\smallskip
\noindent
2) $X$ \underline{posee un $0$-ciclo de grado $1$}.

\smallskip
\noindent
3) \underline{Existe un haz invertible $L$ sobre $X\times P^o$ tal que
el homomorfismo correspondiente}
$P^o_\eta\to\underline{\operatorname{Pic}}_{X_\eta/\eta}=P_\eta$
\underline{sea la inclusión canónica}. (Se sabe, por J.-P. Murre, que
$P_\eta$ es representable por un esquema en grupos localmente de tipo
finito sobre $\eta$; véase, por ejemplo, SGA 6, XII, 1.5.)

%%%%%%%%%%  scan idx 477  |  folio 466  %%%%%%%%%%

\smallskip
\noindent
4) $X$ \underline{es cohomológicamente plano sobre $S$} (EGA III,
7.8.1) (equivalentemente, cohomológicamente plano en dimensión cero, o,
equivalentemente,
$k\xrightarrow{\sim}H^o(X_o,\underline{\mathcal O}_{X_o})$),
\underline{y $d=1$}.

\smallskip
\noindent
5) $P$ \underline{es representable} (equivalentemente, representable
por un esquema en grupos liso sobre $S$).

\smallskip
\noindent
6) $P^o$ \underline{es representable} (equivalentemente, representable
por un esquema en grupos liso sobre $S$). (Para la definición de $P^o$,
véase SGA 3, $\mathrm{VI}_B$, 3.1.)

\smallskip
\noindent
7) $P^o$ \underline{es separado sobre $S$}, es decir, satisface el
criterio valuativo de separación de EGA II, 7.2.3 (equivalentemente,
toda sección $s$ de $P^o$ sobre $S$ tal que $s(\eta)=0$ es nula).

\smallskip
\noindent
8) $d=1$.

\smallskip
\noindent
\underline{Se tienen entonces las implicaciones}

\begin{equation}\tag{12.1.4}
 1)\Longleftrightarrow2)\Longrightarrow3)\Longrightarrow4)
 \Longleftrightarrow5)\Longleftrightarrow6)\Longleftrightarrow7)
 \Longrightarrow8).
\end{equation}

\noindent
\underline{En particular, si $d=d^t$, por ejemplo si $k$ es perfecto,
las condiciones 1) a 8) son equivalentes}.

\medskip
\noindent
b) \underline{Supongamos que $d^t=1$, de modo que $P$ es representable
por un esquema en grupos liso sobre $S$. Entonces la adherencia
esquemática de la sección unidad de $P$ es un subesquema en grupos $E$
étale sobre $S$. Si $I$ es el conjunto de las componentes irreducibles
reducidas $X_o^i$ de $X_o$, se obtiene un isomorfismo}

\begin{equation}\tag{12.1.5}
 \mathbb Z^I\simeq\Gamma(S,E)
\end{equation}

%%%%%%%%%%  scan idx 478  |  folio 467  %%%%%%%%%%

\noindent
\underline{asociando a cada $i\in I$ la sección $S_i$ de $P$ sobre $S$
(nula sobre $\eta$) definida por el haz invertible
$\underline{\mathcal O}_X(X_o^i)$ \sisi{sobre $S$}{sobre $X$}}.
\footnote{La fuente impresa dice aquí «sobre $S$» y escribe después
$X_i^o$ para el mismo divisor. El haz invertible asociado a la componente
$X_o^i$ vive sobre $X$.}
(Como, por hipótesis, los anillos locales de $X$ son factoriales, los
ciclos $X_o^i$ son efectivamente divisores; por tanto,
$\underline{\mathcal O}_X(X_o^i)$ tiene el sentido indicado como haz
invertible asociado a dicho divisor; véase EGA IV, 21.2.8.1.)
\underline{Se tiene}

\begin{equation}\tag{12.1.6}
 P^o\cap E=\underline{\text{la sección nula de}}\ P.
\end{equation}

\noindent
c) \underline{Bajo las condiciones de b), el cociente}

\begin{equation}\tag{12.1.7}
 Q=P/E
\end{equation}

\noindent
\underline{es representable por un esquema en grupos liso y separado
sobre $S$, y este esquema es néroniano, es decir, satisface la propiedad
universal (1.1.2) en términos de su fibra \sisi{especial}{genérica}}.

\medskip
\noindent
d) \underline{Bajo las condiciones de c), supongamos además que la fibra
geométrica genérica de $X$ es íntegra. Entonces existe una sucesión exacta
canónica}

\begin{equation}\tag{12.1.8}
 0\longrightarrow A\longrightarrow Q
   \xrightarrow{\ \deg\ }\mathbb Z_S\longrightarrow0,
\end{equation}

\noindent
\underline{donde $\deg$ está inducido por el homomorfismo «grado»
$P\to\mathbb Z_S$, y donde $A$ es un esquema en grupos liso, separado y
de tipo finito sobre $S$, que también es néroniano, es decir, un modelo
de Néron de}

\begin{equation}\tag{12.1.9}
 A_\eta=P_\eta^o=
 \underline{\operatorname{Pic}}^o_{X_\eta/\eta}.
\end{equation}

\noindent
\underline{Además, el homomorfismo compuesto $P^o\to P\to Q$ induce un
isomorfismo}

\begin{equation}\tag{12.1.10}
 P^o\xrightarrow{\ \sim\ }A^o.
\end{equation}

%%%%%%%%%%  scan idx 479  |  folio 468  %%%%%%%%%%

\medskip
\noindent
12.1.11. El caso que nos interesará es aquel en el que $X_\eta$ es liso
sobre $\eta$, de modo que, teniendo en cuenta (12.1.1), la fibra
geométrica genérica de $X$ es lisa y conexa (por tanto, íntegra), y en el
que $X_\eta$ posee un $0$-ciclo de grado uno. Se cumplen entonces todas
las condiciones consideradas en 12.1 (y, en particular, las condiciones
1) a 8) de a)). Además, $A_\eta=P_\eta^o$ es entonces un esquema
abeliano, y d) nos proporciona una construcción muy explícita de su
modelo de Néron $A$. En particular, (12.1.10) proporciona un isomorfismo

\begin{equation}\tag{12.1.12}
 A_o^{\,o}\simeq\underline{\operatorname{Pic}}^o_{X_o/k},
\end{equation}

\noindent
que muestra, en particular, que $A$ \underline{tiene reducción
semiestable sobre $S$ si y solo si}
$\underline{\operatorname{Pic}}^o_{X_o/k}$ \underline{tiene rango
unipotente cero}. Cuando $X$ es regular y $X_o$ es separable sobre $k$,
se sabe que esto ocurre si y solo si la fibra geométrica especial
$\overline X_o$ no tiene otros puntos singulares que puntos dobles
ordinarios (es decir, puntos cuyo anillo local completado coincide con el
de la unión de los dos ejes coordenados del espacio afín de dimensión
$2$ en el origen). La verificación de este hecho, junto con la
descripción «bien conocida» del esquema de Picard de una curva propia
separable sobre un cuerpo $k$, utilizando EGA IV, 21.8.5, se deja al
lector como ejercicio. Obsérvese también que la implicación
«$\overline X_o$ tiene puntos dobles ordinarios» $\Longrightarrow$
«$A_\eta=\underline{\operatorname{Pic}}^o_{X_\eta/\eta}$ tiene reducción
semiestable» es una consecuencia inmediata del «teorema de monodromía
geométrica» III\sisi{\hspace{3.5em}}{}, teniendo en cuenta el criterio
galoisiano de reducción semiestable de 3.5. Sin embargo, este argumento
no proporciona el importante isomorfismo (12.1.10) ni su consecuencia
(12.1.12).

%%%%%%%%%%  scan idx 480  |  folio 469  %%%%%%%%%%

\medskip
\noindent
12.2. Para la demostración de 12.1 remitimos a [22 bis]; en realidad,
solo necesitaremos el caso en que $X_\eta$ es liso sobre $\eta$. Notemos
también que en [22 bis] se demuestran relaciones entre esquemas de Picard
y modelos de Néron bajo condiciones bastante más generales que las de
12.1 c) y d). Como sus enunciados son más técnicos, hemos preferido no
incluirlos en este recordatorio. Señalemos únicamente que b), c) y d) son
consecuencias fáciles de a), y que en a) los únicos puntos delicados son
las implicaciones $3)\Longrightarrow4)\Longrightarrow5)$ y
$7)\Longrightarrow4)$, siendo la primera la más difícil. Este punto es,
en cambio, trivial si se supone que $X_o$ es separable sobre $k$ (lo que
implica que $X$ admite una sección sobre $S$, pues $S$ es estrictamente
local, y, por tanto, que se cumple la condición 1), $d^t=1$). Aconsejamos
al lector que haga la demostración en este caso particularmente sencillo,
suponiendo además, si es necesario, que $X_\eta$ es liso, de modo que
$P_\eta^o$ sea un esquema abeliano. Se puede utilizar entonces su modelo
de Néron $A$, y $A^o$ representa
$\underline{\operatorname{Pic}}^o_{X/S}$. Solo utilizaremos 12.1 en este
caso particular y bajo la hipótesis adicional de que $\overline X_o$ no
tiene otros puntos singulares que puntos dobles ordinarios, es decir, por
(12.1.11), de que $A_\eta$ tiene reducción semiestable. Este es el caso
considerado en \sisi{12.5.4}{10.5.4}, al cual nos limitaremos de ahora en adelante.
